\documentclass[runningheads]{llncs}
\usepackage[T1]{fontenc}
\usepackage{graphicx}
\usepackage{booktabs}
\usepackage{float}
\usepackage{amsmath}
\usepackage{amssymb}
\usepackage{enumitem}
\usepackage{microtype}
\usepackage{tikz}
\usetikzlibrary{positioning,arrows.meta,calc}
\usepackage{pgfplots}
\pgfplotsset{compat=1.18}
\usepackage{algorithm}
\usepackage[noend]{algpseudocode}
\usepackage[hidelinks]{hyperref}
\setlength{\textfloatsep}{8pt plus 2pt minus 3pt}
\setlength{\floatsep}{6pt plus 1pt minus 2pt}
\setlength{\intextsep}{6pt plus 1pt minus 2pt}
\setlength{\abovecaptionskip}{4pt}
\setlength{\belowcaptionskip}{2pt}
\raggedbottom
\newcommand{\figplaceholder}[2]{%
  \fbox{\parbox[c][#1][c]{0.92\linewidth}{\centering\small\textit{Placeholder -- generate with \texttt{figures/make\_figures.py}}\\[2pt]\footnotesize #2}}}
\begin{document}
\title{Diagnosing the Believability--Reliability Gap in LLM-Agent Simulations of Online Communities}
\titlerunning{Diagnosing the Believability--Reliability Gap}
\author{Anonymous Author(s)}
\authorrunning{Anonymous Author(s)}
\institute{Anonymised for double-blind review}
\maketitle


\begin{abstract}
Generative agents powered by large language models (LLMs) increasingly simulate online
communities for low-cost in silico Web research, and can readily appear convincing: fluent
reasoning, memories, and plans. We ask whether this surface plausibility can be trusted. We
distinguish individual \emph{believability} (how human-plausible an agent's behaviour is) from
collective \emph{task outcome reliability} (whether the simulated community meets the task's
success criterion), scoring both from the same event traces across two online-community
coordination tasks (a commons dilemma, CD, and a public-voting information-consensus task, IC)
under a four-condition architecture comparison plus one exploratory reflection variant (120
primary, 20 exploratory trials). Believability and reliability are positively linked but the link
is fragile (Spearman $\rho=0.55$ in CD, $0.75$ in IC). It is carried almost entirely
by \emph{which cognitive modules an architecture has} rather than how well an agent uses them, and
it disappears or reverses when modules an architecture lacks are scored as not applicable rather
than zero. Believable agents are therefore neither necessary nor sufficient for reliable outcomes. The
sharpest evidence is the voting task: every agent reasons plausibly yet the group still fails,
abandoning correct private information once the public tally appears and converging unanimously
on the wrong answer whenever the correct option is the salient first-listed one. Counterbalancing
shows this is confounded with option position, so we describe it as cascade-like herding rather
than pure social-information lock-in. A coordination gain is \emph{associated with} an added
reflection rule, but a matched content-free placebo recovers a substantial portion of the gain, so
it is not clearly attributable to the rule's content. The practical warning: a believable
simulation is not yet a reliable one, and the two must be measured separately.
\keywords{Generative agents \and LLM multi-agent systems \and Agent-based simulation \and Simulation validation \and Herding and social-information lock-in \and Online communities \and Web analytics.}
\end{abstract}


\section{Introduction}\label{sec:intro}
Web research increasingly uses \emph{generative agents} (LLM-driven participants conditioned on
a persona, memories, goals, and history) in place of hand-coded
rules~\cite{park2023generative,gao2024survey,ghaffarzadegan2024tutorial}, with simulators now
running up to a million agents~\cite{yang2024oasis,piao2025agentsociety}. Because these agents
look convincing, it is tempting to treat plausibility as evidence the simulation is sound.
Agent-based modelling warns against this: validity requires verification and validation, not
appearances~\cite{windrum2007empirical,moss2005towards}. We therefore separate
\emph{believability}, a \emph{micro} property (how human-plausible a single agent's behaviour is),
from \emph{collective task outcome reliability}, a \emph{macro} property (whether the community
meets the task's success criterion). A simulation can fail yet look realistic, or succeed for
unrealistic reasons, so we make no claim about correspondence to a real system. The open question,
not previously measured directly, is their within-simulation relationship: does individual
believability imply collective reliability?

This matters for the Web because online mechanisms (ratings, voting, trending lists, comment
threads) expose earlier choices to later participants, so reasonable individual behaviour can
produce a collectively wrong \emph{information cascade}: agents discarding private information to
follow the visible majority~\cite{bikhchandani1992theory}, a pattern LLM agents
reproduce~\cite{cho2025herd,jain2025social,zhong2025conformity}. A simulation can thus be full of
believable agents while the community converges confidently on the wrong answer. We test this in
two structurally distinct coordination tasks (a commons dilemma and a public-voting
information-consensus task) under a \emph{cumulative} architecture ladder (baseline;
$+$memory; $+$planning; $+$reflection) plus one exploratory LLM-generated-reflection variant (120
primary, 20 exploratory trials). Being cumulative, the ladder is a capability comparison, not a
clean ablation. We contribute:
\begin{enumerate}
  \item \textbf{A dual-level measurement instrument} scoring individual believability $B$ (four
  automatic sub-dimensions, two via an LLM-as-judge rubric~\cite{zheng2023judging}) and collective
  task outcome reliability $V$ from the same traces, plus a blinded human-rating protocol
  ($\alpha\geq0.67$ reliability gate) and a three-rater pilot sanity check.
  \item \textbf{A measurable believability--outcome gap}, summarised by a signed discrepancy
  $\Delta=B-V$ and a two-dimensional $(B,V)$ view. The strongest evidence is the voting task:
  believable agents converge unanimously onto the wrong answer when it is not the salient
  first-listed option. Counterbalancing shows this is cascade-like herding confounded with
  position rather than pure lock-in.
  \item \textbf{An account of where reliable outcomes come from:} the link is associated with which
  modules an architecture \emph{has} rather than how well agents express them, disappearing or
  reversing under not-applicable scoring, with no believability threshold for success. Separately,
  a categorical coordination jump ($10\%\rightarrow70\%$) is associated with a
  deterministic reflection rule. A matched content-free placebo reproduces most of the gain, though,
  so we treat the link as associative rather than causal~\cite{aharon2026tacit}.
  \item \textbf{A failure-trace taxonomy} distinguishing two system-level failure mechanisms from
  logged state alone: persistent oversubscription versus sub-threshold believability--outcome
  failure.
\end{enumerate}
\noindent An anonymised replication package (code, run manifests, raw event logs, and
analysis scripts) accompanies this submission.\footnote{\url{https://anonymous.4open.science/r/believability-reliability-gap}}


\section{Related Work}\label{sec:related}
\textbf{Validation in (generative) ABM.} Agent-based modelling studies how macro patterns emerge from interacting agents~\cite{epstein1996growing}, with a mature validation vocabulary: \emph{verification} (built correctly?) vs.\ \emph{validation} (right model?)~\cite{sargent2010verification,rand2011rigor}, plus conceptual, operational, and data validity~\cite{sargent2010verification,augusiak2014evaludation}, standard instruments (ODD~\cite{grimm2006odd,grimm2020odd}, pattern-oriented modelling~\cite{grimm2005pom}), generative sufficiency~\cite{moss2005towards}, mechanism plausibility~\cite{zhao2026mechanism}, and the \emph{micro-face} vs.\ \emph{macro-face} distinction~\cite{windrum2007empirical}. Kl\"ugl identifies the central difficulty as the \emph{missing link between a model's micro and macro levels}~\cite{kluegl2008validation}, which is exactly where our study sits~\cite{augusiak2014evaludation}. For GABMs, recent work finds validation largely unmet~\cite{larooij2025critical}, an ``uncanny valley'' substituting fluent output for emergence~\cite{zeng2026uncanny}, and under-reported variance~\cite{wu2026boundary}. Proposed responses (validation workflows~\cite{adornetto2025overview}, replication~\cite{cross2025validating}, realism checks~\cite{munker2026benchmark}, platform calibration~\cite{tomasevic2026operational}) target external fidelity rather than the micro--macro relationship we study.

\textbf{LLM-agent evaluation and coordination.} LLM-agent evaluation has matured around capability benchmarks testing agents in isolation~\cite{li2024agentbench,mohammadi2025evaluation}, but multi-agent settings surface failures invisible individually: a communicating--reasoning gap~\cite{zhang2026silo}, destabilisation by one uncooperative agent~\cite{kulshreshtha2026defection}, and memory-design effects on cooperation~\cite{hishiki2026memory}. Two mechanisms matter here: \emph{informational herding} (discounting private signals for the public majority~\cite{cho2025herd,jain2025social,zhong2025conformity}, echoing information cascades~\cite{bikhchandani1992theory}) and \emph{tacit focal-point coordination} (coordinating on a salient cue rather than reasoning~\cite{aharon2026tacit}). These underlie, respectively, our clearest believable-but-unreliable case and our largest architectural effect.

\textbf{Trace-based evaluation.} Treating logs as event traces links GABM validation to predictive process monitoring~\cite{galanti2020explainable}: a single end-of-run score cannot explain \emph{why} a run failed~\cite{cemri2025fails}, and system-, trace-, and node-level evaluation has been systematised~\cite{yehudai2026clear}. Since automated judges can diverge from humans~\cite{pan2025agentrewardbench}, we pair an LLM judge~\cite{zheng2023judging} with a human-validation gate, adding two simulation-specific failure types (D, E). Table~\ref{tab:related} positions our study on six dimensions; we are aware of no prior study that jointly measures micro believability and macro outcomes \emph{at the trial level}.
\begin{table}[t]
\centering
\caption{Where prior work sits on six evaluation dimensions ($\bullet$ = addressed, $\circ$ = partial): Micro/Macro = individual/collective evaluation; Calib.\ = empirical calibration; Mech.\ = mechanism analysis; Trace = trace-based diagnostics; Human = human validation (This work: pilot sanity check only).}
\label{tab:related}
\small
\setlength{\tabcolsep}{4.5pt}
\begin{tabular}{lcccccc}
\toprule
Work & Micro & Macro & Calib. & Mech. & Trace & Human \\
\midrule
GABM critical review~\cite{larooij2025critical}      & $\circ$ & $\circ$ & $\circ$ & $\circ$ & --      & -- \\
Operational calibration~\cite{tomasevic2026operational} & --      & $\bullet$ & $\bullet$ & $\circ$ & --   & -- \\
ODD / pattern-oriented~\cite{grimm2020odd,grimm2005pom} & --   & $\bullet$ & $\bullet$ & $\circ$ & --   & -- \\
Multi-agent failure taxonomy~\cite{cemri2025fails}   & --      & $\circ$ & --      & $\bullet$ & $\bullet$ & $\circ$ \\
Agent eval surveys~\cite{mohammadi2025evaluation,yehudai2026clear} & $\bullet$ & $\circ$ & -- & $\circ$ & $\bullet$ & $\circ$ \\
\textbf{This work} & $\bullet$ & $\bullet$ & -- & $\bullet$ & $\bullet$ & $\circ$ \\
\bottomrule
\end{tabular}
\end{table}


\section{A Dual-Level Validation Instrument}\label{sec:instrument}

\subsection{Preliminaries}\label{sec:prelim}
We study generative multi-agent simulations at two levels. A \emph{trial} $t$ is one run of a
coordination task by a fixed group of $n$ agents under an architecture \emph{condition}
$\kappa\in\{C_1,\dots,C_5\}$ that varies the agents' cognitive modules; it proceeds in discrete
\emph{steps}, and at each step every agent emits a structured action, yielding the event traces
from which all measures are computed. The instrument assigns each trial two scalars in $[0,1]$:
\textbf{individual believability} $B(t)$, how human-plausible the agents' behaviour is at the
\emph{micro} level (given persona, memory, and goals), and \textbf{collective task outcome
reliability} $V(t)$, how well the group meets the task's success criterion at the \emph{macro}
level. A trial \emph{succeeds} when $V(t)\geq\theta$ for a task-specific threshold $\theta$.
Throughout, ``believability'' is micro-level plausibility only, with no claim about
correspondence to a real system; $B$ and $V$ are both internal, within-simulation measures.

\subsection{Problem Statement}\label{sec:problem}
The problem we address is whether believability $B$ implies collective task outcome reliability
$V$: whether agents that look believable in isolation produce a community that reliably succeeds,
or whether believable agents can coordinate into collectively unreliable outcomes. Current practice
implicitly assumes high believability is sufficient for valid collective behaviour, i.e.\ that a
single \emph{believability gate} exists: a threshold $\tau$ such that, for every trial,
\begin{equation}\label{eq:gate}
B(t)\;\geq\;\tau \;\;\Longrightarrow\;\; V(t)\;\geq\;\theta .
\end{equation}
For this to be a non-vacuous, practically useful claim, $\tau$ must be exceeded by at least some
observed trials; a $\tau$ above every observed $B(t)$ satisfies Eq.~\eqref{eq:gate} trivially and
is not informative. To examine the two scalars jointly we use their signed difference, the
\emph{believability--outcome discrepancy}
\begin{equation}\label{eq:discrepancy}
\Delta(t)\;=\;B(t)-V(t).
\end{equation}
A \emph{gap case} is a trial where believable agents nonetheless fail the task: a high discrepancy
\emph{and} an outcome below threshold, collected into
\begin{equation}\label{eq:gapset}
\mathcal{G}\;=\;\{\, t : \Delta(t) > \delta \ \text{and} \ V(t) < \theta \,\}.
\end{equation}
The question is therefore whether some non-vacuous gate $\tau$ makes the gap set empty
($\mathcal{G}=\varnothing$), or whether believable agents can be collectively unreliable
($\mathcal{G}\neq\varnothing$, with no such $\tau$).



\subsection{Micro Level: Individual Believability}\label{sec:micro}
Believability $B$ is a weighted composite of four automatically-scored sub-dimensions
$b_k\in[0,1]$ with weights $w_k$, fixed a priori rather than empirically tuned:
\begin{equation}\label{eq:believability}
B \;=\; \sum_{k=1}^{4} w_k\, b_k .
\end{equation}
\textbf{Behavioural consistency} $b_1=\sum_{j=1}^{3}\gamma_j c_j$ combines request consistency,
persona alignment (cosine similarity via OpenAI's \texttt{text-embedding-3-small}), and
cooperation rate (also reported with it removed, since it overlaps the macro outcome).
\textbf{Memory coherence} $b_2$ combines citation frequency with the cosine relevance of cited
memories. \textbf{Planning plausibility} $b_3$ and \textbf{response naturalness} $b_4$ use an
LLM-as-judge rubric~\cite{zheng2023judging,pan2025agentrewardbench}, blind to condition, outcome,
and architecture; the judge is \texttt{gpt-4o-mini} (the generator's model), so these two
sub-dimensions are not judge-independent (revisited in the limitations); we emphasize categorical
outcomes and report robustness checks for the claims that depend on $B$. A cross-provider judge is
available in the replication package but was not run for the reported results.

\textbf{Capability presence versus available-dimension scoring.} A module an architecture lacks
(e.g.\ memory in the baseline) cannot be expressed, so scoring it $0$ conflates \emph{absence} with
\emph{poor use}. We therefore separate capability \emph{presence}, \emph{quality of use}, and
\emph{architecture-independent} believability (consistency and naturalness, expressible
everywhere), reporting the composite two ways: presence-inclusive (absent dimensions $=0$) and
not-applicable (weights renormalised over available dimensions).

\textbf{Instrument validation.} A blinded human-rating protocol scores condition-blinded decision
packets, balanced over conditions, outcomes, phases, and personas, on anchored $1$--$5$ rubrics;
inter-rater reliability uses Krippendorff's $\alpha$. A three-rater, $15$-packet pilot is a
\emph{sanity check} rather than convergent validity: human and automated composites rank together at
$\rho=0.82$ ($n=15$), carried almost entirely by planning ($\rho=0.79$), and no dimension reaches
the $\alpha\geq0.67$ gate, so full multi-rater validation remains outstanding.




\subsection{Macro Level: Collective Task Outcome Reliability}\label{sec:macro}
Collective task outcome reliability $V$ is computed from observable simulation state, so it is
reproducible, but measures task outcomes rather than validity in the modelling sense. Its metrics are
coordination success, \textbf{coordination score} (fraction of coordinated steps),
\textbf{sustainability} (pool level over its initial level), \textbf{demand pressure}, the
\textbf{Gini coefficient} over \emph{granted} allocations, \textbf{convergence speed}, and
\textbf{failure traceability} against a taxonomy of agent-level causes~\cite{cemri2025fails}.

For the commons task, $V$ is a weighted sum of three outcome terms $v_m$ (coordination score
$v_1$, sustainability $v_2$, and equity $v_3 = 1 - \text{Gini}$) with weights $\omega_m$ fixed a
priori (coordination weighted above sustainability and equity; robustness across a range of
$\omega$ is reported with the gap-case analysis):
\begin{equation}\label{eq:vscore}
V_{\mathrm{CD}} \;=\; \sum_{m=1}^{3} \omega_m\, v_m .
\end{equation}
For the consensus task, $V$ is the coordination score $v_1$ alone. Both forms lie in $[0,1]$, and a
trial succeeds when $V \geq \theta$ for the task's success threshold $\theta$.

\textbf{Failure taxonomy.} Beyond the coordination, specification, and verification failures
(Types~A--C) of prior work~\cite{cemri2025fails}, we add two system-level types: \emph{Type~D},
persistent oversubscription (demand exceeds replenishment), and \emph{Type~E}, sub-threshold
failure where agents are individually believable yet the collective outcome falls short: the
believability--outcome gap case. Both are assigned automatically from logged state. Because
Type~E is defined partly through the discrepancy, we treat it as a diagnostic \emph{label} rather
than independent evidence.


\subsection{Confronting the Two Levels}\label{sec:confront}
Each trial yields a pair $(B(t),V(t))$ in the unit square, the \emph{believability--reliability
plane} ($B$ horizontal, $V$ vertical), our primary object of analysis. We use $\Delta(t)=B(t)-V(t)$
only as a one-dimensional summary for counting and ranking, rather than as a calibrated quantity,
since subtracting $V$ from $B$ presumes a scale comparability we do not assume.

The plane makes the question geometric: a believability gate is a \emph{vertical} line $B=\tau$
right of which every trial also clears $V=\theta$:
\[
\{\,t : B(t)\geq\tau\,\} \;\subseteq\; \{\,t : V(t)\geq\theta\,\}.
\]
Gap cases ($\Delta(t)>\delta$: high $B$ but $V(t)<\theta$) occupy the lower-right; any trial there
means a non-empty $\mathcal{G}$, which rules out a gate $\tau$. Since $\delta$ is a convention
rather than a calibrated boundary, we report how $|\mathcal{G}|$ changes as $\delta$ and $\omega_m$ vary,
treating as informative only gap cases persisting across that range. The same traces also support
\emph{trace-based diagnostics}: why a believable trial failed and which failure type (A--E)
applies.




\section{Experimental Design}\label{sec:design}

\subsection{Agent Architecture and Conditions}
Agents follow Park et al.'s generative-agent design~\cite{park2023generative} (persona, memory,
planning, reflection). A fixed panel of $n=8$ personas, drawn once (seed $42$) from a 25-persona
environment, is held constant across \emph{all} trials, which differ only in per-trial seed and
LLM stochasticity; this shared panel induces a cross-trial dependence we account for statistically.
GPT-4o-mini (\texttt{gpt-4o-mini-2024-07-18}, temperature $0.7$) generates all behaviour; exact
prompts, limits, and run dates are in the replication package. Conditions form a cumulative
ladder: \textbf{C1} baseline; \textbf{C2} $+$memory; \textbf{C3} $+$planning; \textbf{C4}
$+$\emph{deterministic} reflection (a fixed rule appending a fair-share statement each round,
which risks reading as solution injection rather than reflection -- hence the placebo-reflection
control below); \textbf{C5} (exploratory) replaces C4's rule with LLM-\emph{generated} reflection.
C1--C4 are primary. Being cumulative, the ladder does not isolate any single module; we leave a
planning-only condition to future work.




\subsection{Task 1: Commons Dilemma (CD)}
The $n=8$ agents share a renewable community fund, a deliberately demanding \emph{stress-test
abstraction} of Web resource governance (mutual-aid groups, open-source maintainer pools, online
commons), not a calibrated model. Defaults: initial pool 1000 credits, replenishment 50/round,
per-agent request cap 100, capacity 1000; the cap is deliberately far above fair share ($6.25$) so
that over-requesting is a genuine behavioural choice, not an artefact of a narrow action space.
Each step, agents receive a plain-English resource
context and emit a structured request; Algorithm~\ref{alg:env} (left) gives the update. A trial
runs 100 steps unless the pool collapses, and succeeds if there is no collapse and the coordination
score reaches $\theta_{\mathrm{CD}}=0.7$. Episodic memory (C2--C5) records prior rounds (demand
vs.\ replenishment, fair-share deviations, over-/under-requesters). Primary data: $4\times20=80$
trials, plus 10 exploratory C5.



\begin{algorithm}[t]
\caption{Per-step environment update. \emph{Left (CD):} proportional rationing when demand exceeds the pool (any floor remainder stays in the pool for the next round); collapse ends the trial. \emph{Right (IC):} agents vote after seeing the tally; success requires a five-step streak of $\geq\lceil 0.75n\rceil$ correct votes.}
\label{alg:env}
\footnotesize
\setlength{\tabcolsep}{5pt}
\begin{tabular}{@{}p{0.46\linewidth}|p{0.46\linewidth}@{}}
\centering\textbf{\footnotesize Commons dilemma (CD)}
&
\centering\textbf{\footnotesize Information consensus (IC)}
\tabularnewline[1pt]
\begin{minipage}[t]{\linewidth}
\begin{algorithmic}[1]
\Require pool $r$, replenishment $\rho$, capacity $\kappa$, agents $1,\dots,n$
\State $\textit{fair} \gets \rho / n$ \Comment{$=6.25$}
\State broadcast $(r,\ r/r_0,\ \rho,\ n,\ \textit{fair})$
\State collect request $q_i$ capped at $100$
\State $Q \gets \sum_i q_i$ \Comment{total demand}
\If{$Q \le r$}
\State $g_i \gets q_i$ \Comment{full grant}
\Else
\State $g_i \gets \lfloor q_i \cdot r / Q \rfloor$ \Comment{rationing}
\EndIf
\State $r \gets \min(\kappa,\ r - \sum_i g_i + \rho)$
\If{$r \le 0$}
\State \textbf{collapse}; end trial
\EndIf
\State $\textit{coordinated} \gets [\,Q \le \rho\,]$
\State record $r/r_0$, $\mathrm{Gini}(g)$
\end{algorithmic}
\end{minipage}
&
\begin{minipage}[t]{\linewidth}
\begin{algorithmic}[1]
\Require correct option $a^\star$, signals $(A{:}4,B{:}2,C{:}2)$, agents $1,\dots,n$, streak $s \gets 0$
\State broadcast previous public tally $T$
\State collect vote $v_i$ from each agent $i$
\State $T \gets$ tally of $\{v_i\}$ over options
\State $c \gets |\{i : v_i = a^\star\}|$ \Comment{\# correct}
\State $\textit{coordinated} \gets [\,c \ge \lceil 0.75n\rceil\,]$
\If{\textit{coordinated}}
\State $s \gets s+1$
\Else
\State $s \gets 0$
\EndIf
\If{$s \ge 5$}
\State \textbf{success}; record convergence step
\EndIf
\State record $c/n$, $\mathrm{Gini}(T)$
\end{algorithmic}
\end{minipage}
\end{tabular}
\end{algorithm}



\subsection{Task 2: Repeated Public-Voting / Information Consensus (IC)}
The $n=8$ agents must collectively identify which of three options is correct. Private signals give
the correct option a strict plurality (four agents receive the correct signal, two each the
alternatives), so honestly pooling the signals would identify it; each agent sees both its own
private signal and the running public tally. Primary trials fix Option A (first-listed) as
correct; §5.2 counterbalances this choice. This is a direct abstraction of repeated public voting on the
Web (ratings, polls, upvote threads), where herding and social-information lock-in can
arise~\cite{bikhchandani1992theory}; being \emph{repeated simultaneous} voting with a visible tally
rather than the canonical \emph{sequential} cascade, we call the outcomes cascade-like. A step is
coordinated when a fraction $\theta_{\mathrm{IC}}=0.75$ (six of eight) vote correctly, and a trial
succeeds on five consecutive coordinated steps (the convergence step starts that streak);
Algorithm~\ref{alg:env} (right) gives the update. Primary trials \emph{fix} the option set and
correct option; we probe this with counterbalancing and neutral labels below, and leave
hidden-tally and vote-only controls to future work. Primary data: $4\times10=40$ trials, plus 10
exploratory C5.


\subsection{Logging and Analysis}
Composite weights were fixed a priori, ordered by how directly each term reflects its construct:
believability weights $w=(0.30,0.25,0.25,0.20)$ over $(b_1,\dots,b_4)$ (consistency highest, the
two cognitive capabilities equal, naturalness lowest); consistency sub-weights
$\gamma=(0.45,0.35,0.20)$ (cooperation rate least, since it overlaps the macro outcome); commons
outcome weights $\omega=(0.5,0.3,0.2)$ over coordination, sustainability, equity. Being
judgement-based, they carry no conclusion; the gap-case analysis reports robustness across a range
of $\omega$.

Each trial emits micro (per agent--step) and macro (pool, totals, sustainability, inequality,
collapse) logs. Hypotheses were fixed in writing and timestamped before analysis, pre-stated but
not formally preregistered, separating confirmatory analyses (pre-stated associations and
condition comparisons) from exploratory ones (sub-dimension drivers, all C5). We use Spearman
correlation and Mann--Whitney $U$ for association; Kruskal--Wallis and pairwise Mann--Whitney with
Holm correction and Cliff's $\delta$; Wilson intervals for success rates; OLS with standardised
coefficients; and string deduplication for reflection repetition. Because trials share one persona
panel, effective sample size is below nominal and $p$-values are optimistic.

\subsection{Robustness Controls}\label{sec:controls}
Four further experiments probe the main threats (Table~\ref{tab:controls}), all reusing the same
generator, personas, and logging.
\begin{table}[t]
\centering
\caption{Robustness controls. All reuse the same generator, personas, and logging as the main runs.}
\label{tab:controls}
\small
\setlength{\tabcolsep}{4pt}
\begin{tabular}{@{}p{2.1cm} p{2.0cm} p{6.4cm} r@{}}
\toprule
Control & Separates & Manipulation & $n$ \\
\midrule
Placebo reflection (CD) & content from step & injection: fair-share string verbatim; placebo: matched string, no content & 10/arm \\
IC counterbalancing & lock-in from salience & correct option rotated A/B/C; also neutral labels & 5/cell \\
No-plurality-hint memory (IC) & memory from scaffold & plurality hint removed; harder $3{:}3{:}2$ split & 5/cell \\
Persona panels (CD) & single-panel effect & baseline \& full on 3 further panels (seeds $43$--$45$) & 5/cell \\
\bottomrule
\end{tabular}
\end{table}
The persona-panel control gives four panels in total for a mixed-effects analysis (condition
fixed, panel random); a fifth panel (seed 46) was left incomplete and is excluded.



\section{Results}\label{sec:results}
Table~\ref{tab:results} shows the central pattern: presence-inclusive believability rises with
architecture but is nearly flat once memory and planning are present (C3$\to$C4: $+0.013$), while
collective outcomes move categorically, and in the consensus baseline the two visibly part company.
Categorical outcomes are computed from objective simulation state, independent of the believability
instrument; only the association and predictor analyses depend on it, where
zero-versus-not-applicable scoring matters most.
\begin{table}[t]
\centering
\caption{Condition-level results. CD: commons dilemma (20 trials/condition); IC: repeated
public-voting (10 trials/condition). C5 is exploratory (10 trials each). Believ.\ = mean
presence-inclusive composite believability; Sust.\ = sustainability; Succ.\ = coordination
success (CD: Wilson 95\% CI in brackets); Conv.\ = convergence step.}
\label{tab:results}
\small
\setlength{\tabcolsep}{4pt}
\begin{tabular}{l rrr c rrr}
\toprule
 & \multicolumn{3}{c}{CD} & & \multicolumn{3}{c}{IC} \\
\cmidrule{2-4}\cmidrule{6-8}
Condition & Believ. & Sust. & Succ. & & Believ. & Succ. & Conv. \\
\midrule
C1 Baseline           & 0.347 & 0.138 & 0\% [0,16]   & & 0.385 & \textbf{0\%} & timeout \\
C2 Memory             & 0.543 & 0.264 & 5\% [1,24]   & & 0.561 & 100\% & 1.0 \\
C3 Memory + Planning  & 0.669 & 0.326 & 10\% [3,30]  & & 0.669 & 100\% & 1.0 \\
C4 Full (deterministic) & 0.682 & 0.931 & \textbf{70\% [48,85]} & & 0.668 & 100\% & 1.0 \\
\midrule
\emph{C5 Full (LLM reflection, exploratory)} & \emph{0.674} & \emph{0.988} & \emph{100\%} & & \emph{0.669} & \emph{100\%} & \emph{1.1} \\
\bottomrule
\end{tabular}
\end{table}



\subsection{A Fragile, Presence-Driven Association; No Believability Gate}\label{sec:results-assoc}
With the presence-inclusive composite, believability is positively associated with coordination, but the association is \emph{fragile} (Table~\ref{tab:assoc}): it survives a persona-clustered bootstrap and replicates across three further panels (strong mixed-effects condition effect, negligible between-panel variance), yet falls to $\rho=-0.24$ under not-applicable scoring and $\rho=0.23$ on the architecture-independent composite, so it is carried by capability \emph{presence} rather than expressed believability. There is also \emph{no believability threshold} for success: the most-believable failing trial ($0.698$) exceeds the least-believable succeeding one ($0.658$), and the floor near $0.36$--$0.39$ is a presence artefact, not a behavioural gate. Within these tasks, no believability threshold is necessary or sufficient for reliable outcomes (Fig.~\ref{fig:scatter}).
\begin{table}[t]
\centering
\caption{Robustness of the believability--coordination association. Spearman $\rho$ (CD $n=80$, IC $n=40$; 95\% CIs in brackets where computed). The headline row is trial-level; the persona-clustered bootstrap instead resamples personas and computes $\rho$ at the agent level, a different unit, so its interval need not bracket the headline estimate, and IC's near-zero width reflects its enriched conditions being near-deterministic (100\% success). ``--'' = not informative (IC enriched conditions are degenerate).}
\label{tab:assoc}
\small
\setlength{\tabcolsep}{5pt}
\begin{tabular}{l cc}
\toprule
Scoring / robustness variant & CD $\rho$ & IC $\rho$ \\
\midrule
Presence-inclusive composite (headline, trial-level) & 0.55 [0.38, 0.69] & 0.75 [0.57, 0.85] \\
Persona-clustered bootstrap (agent-level $\rho$) & [0.45, 0.52] & [0.75, 0.75] \\
Not-applicable scoring & $-0.24$ & -- \\
Architecture-independent (consistency $+$ naturalness) & 0.23 & -- \\
Cooperation-rate component removed & 0.44 & -- \\
\midrule
Persona panels 43--45 (mixed-effects, full vs.\ baseline) & \multicolumn{2}{c}{$\beta=0.74$, $p<0.001$, ICC $\approx 0$} \\
\bottomrule
\end{tabular}
\end{table}
\begin{figure}[t]
\centering
\IfFileExists{figures/fig_scatter.pdf}%
  {\includegraphics[width=0.68\textwidth]{figures/fig_scatter.pdf}}%
  {\figplaceholder{3.9cm}{Trial-level $(B,V)$ scatter, split into CD/IC panels: composite believability ($x$) vs.\ collective task outcome $V$ ($y$), colour by condition; gap cases ($\Delta>0.20$) ring-highlighted.}}
\caption{The two-dimensional $(B,V)$ view, primary conditions C1--C4 only (C5 is exploratory and omitted). Horizontal axis: composite believability $B$ (presence-inclusive); vertical axis: collective task outcome reliability $V$; dashed line: success threshold $V=\theta$ (CD $0.70$, IC $0.75$). Ringed points are gap cases ($\Delta > 0.20$).}
\label{fig:scatter}
\end{figure}




\subsection{The Believability--Outcome Gap}\label{sec:results-gap}
The positive association does not imply reliable outcomes (our central result). Per-trial $\Delta = B - V$ flags 51 of 120 primary trials as gap cases at $\Delta>0.20$ (41 CD, all 10 IC baseline); the CD count is choice-sensitive ($37$--$56$ across thresholds; $30$--$59$ across $V$-weightings; bootstrap $[32,50]$), whereas the IC count is invariant ($V=0$, $B$ above the margin) and coded Type~E. The qualitative point is robust: believability does not deliver reliable outcomes on its own.

The IC baseline is the sharpest, instrument-independent demonstration (Fig.~\ref{fig:cascade}): agent reasoning is scored as highly plausible (consistency $0.69$, naturalness $0.89$), yet all 10 trials converge \emph{unanimously on an incorrect option} (by step~2 in 7 of 10), and all 80 first-round votes follow the agent's \emph{own} private signal. Abandonment begins only after the tally is seen. Counterbalancing shows the failure is confounded with option \emph{position} (Table~\ref{tab:counterbal}): the baseline fails only when the correct option is first-listed, under both meaningful and neutral labels. We therefore read it as cascade-like herding~\cite{bikhchandani1992theory,cho2025herd,aharon2026tacit} rather than pure signal-discounting lock-in, a regime that face-validity practice~\cite{kluegl2008validation,larooij2025critical,zeng2026uncanny} is poorly placed to detect. Enriched conditions (C2--C5) all reach 100\%; a no-flag control (hint removed, a harder $3{:}3{:}2$ split in which the correct option is tied with an alternative rather than holding a plurality) still lets memory converge (baseline $0/5$, memory $4/5$), pointing to genuine aggregation rather than a scaffold. Fine-grained statistics therefore draw on CD.
\begin{table}[t]
\centering
\caption{IC option counterbalancing (5 trials per cell): coordination success by which option is correct, under meaningful (A/B/C) and neutral (X/Y) labels.}
\label{tab:counterbal}
\small
\setlength{\tabcolsep}{6pt}
\begin{tabular}{l ccc c cc}
\toprule
 & \multicolumn{3}{c}{Meaningful labels} & & \multicolumn{2}{c}{Neutral labels} \\
\cmidrule{2-4}\cmidrule{6-7}
Condition & A (1st) & B (2nd) & C (3rd) & & X (1st) & Y (2nd) \\
\midrule
Baseline & 20\% & 100\% & 100\% & & 0\% & 100\% \\
Full & 100\% & 100\% & 80\% & & 100\% & 100\% \\
\bottomrule
\end{tabular}
\end{table}
A plausible mechanism, visible in the traces: the correct option's signal support (4/8) is a
plurality, not a majority, and agents reading the tally appear to weigh the \emph{combined} support
for the remaining options against their own single-option read, tipping toward whichever
alternative accumulates votes first (Fig.~\ref{fig:cascade}). When the correct option is not
first-listed, the same trailing dynamic instead works against the wrong options. Neutral labels
do not remove the effect (Table~\ref{tab:counterbal}, X/Y), ruling out semantic association with
``A'' specifically and pointing to a structural comparison bias in how the tally is aggregated.
\begin{figure}[t]
\centering
\IfFileExists{figures/fig_cascade.pdf}%
  {\includegraphics[width=0.68\textwidth]{figures/fig_cascade.pdf}}%
  {\figplaceholder{3.6cm}{Vote-tally trajectories of one representative IC baseline trial.}}
\vspace{-2.5mm}
\caption{Anatomy of a believable-but-unreliable trial (IC baseline, trial 0): the correct option starts with a 4/8 signal plurality, yet the tally tips to an alternative and locks in unanimously by step 2. One agent's flip (excerpted, with a bracketed clarification): ``\emph{While my private signal points to Option A, the combined support for [the alternatives] suggests a stronger collective inclination\,\ldots}'', individually plausible yet collectively wrong.}
\label{fig:cascade}
\vspace{-4mm}
\end{figure}




\subsection{Capability Presence, Not Expressed Quality, Carries the Association}\label{sec:results-drivers}
Separating capability presence, quality of use, and architecture-independent believability (CD, $n=80$; IC enriched conditions are degenerate) isolates the correlates of the association (Table~\ref{tab:drivers}): presence is the strongest signal, quality of use predicts coordination only for memory, and architecture-independent believability is weak. A joint standardised OLS finds planning significant ($\beta=0.45$, $p=0.0012$) but not memory ($p=0.556$), contrary to our pre-stated expectation~\cite{park2023generative,hishiki2026memory}. Within a fixed architecture these associations are unstable at this sample size (planning $\rho=0.13$, n.s.; within-C4 CI $[-0.10,0.82]$), reflecting low power. In short: capability \emph{presence} moves both believability and outcomes, while \emph{expressed} quality moves believability far more than outcomes (the only robust pairwise gain, memory\_planning$\to$full in Table~\ref{tab:results}, is into reflection: Cliff's $\delta\geq0.79$, Holm $p<0.0001$). For builders, sub-dimension metrics best indicate \emph{what an architecture has} rather than how reliably it will coordinate.
\begin{table}[htbp]
\centering
\caption{Sub-dimension drivers of coordination (CD primary, $n=80$; Spearman $\rho$). Presence indicators dominate; expressed quality matters only for memory. n.s.\ = not significant.}
\label{tab:drivers}
\small
\setlength{\tabcolsep}{6pt}
\begin{tabular}{l l c c}
\toprule
Group & Predictor & $\rho$ & $p$ \\
\midrule
Capability presence & planning-present & 0.51 & $<0.001$ \\
                    & memory-present & 0.35 & 0.0015 \\
Quality of use & memory coherence (memory-bearing) & 0.46 & 0.0002 \\
               & planning plausibility (planning-bearing) & 0.13 & 0.44 (n.s.) \\
Architecture-independent & consistency & 0.25 & \\
                        & naturalness & $-0.11$ & n.s. \\
\midrule
\multicolumn{4}{l}{Joint standardised OLS: planning significant ($\beta=0.45$, $p=0.0012$); memory not ($p=0.556$).} \\
\bottomrule
\end{tabular}
\end{table}


\subsection{Coordination Gain Associated with Repetitive Normative Reflection}\label{sec:results-reflection}
C4 (adding deterministic reflection to C3) shows the only robust categorical gain, CD coordination $10\%\rightarrow70\%$ (sustainability $0.326\rightarrow0.931$), yet the reflection signal is highly repetitive (Table~\ref{tab:refl}, Fig.~\ref{fig:saturation}): three unique strings across all 47{,}040 C4 instances, 95.8\% covered by two, saturating by step~2. Because that norm is close to stating the solution, C4 alone cannot separate a direct-instruction effect from a coordination effect. A matched placebo tests this: a tone-, length-, and register-matched string with \emph{no} coordination content shows no statistically significant difference from the fair-share injection (Mann--Whitney $U$ test $p=0.12$, $\delta=0.42$), and both lift coordination far above C3's $10\%$, so most of the gain comes from the reflection \emph{step}, not the \emph{content}. C5 agrees: free LLM reflection yields 1{,}190 unique strings yet 94.9\% restate the same norm (classified by fair-share/replenishment/sustainability keyword matching) and also reaches 100\%. The safer reading: the condition works by repeatedly injecting a shared coordination norm; focal-point coordination~\cite{aharon2026tacit}, direct instruction, and solution leakage remain unseparated. We thus report a gain \emph{associated with} the condition.
\begin{table}[t]
\centering
\caption{Reflection analysis (commons dilemma). String counts are exact-match after normalisation; ``Success'' is coordination success. Instance counts are per-step reflection-context snapshots (8 agents $\times$ 100 steps/trial, growing 0$\to$1$\to$2 then holding at 3 once saturated) summed over 20 trials for C4 and 10 for C5, hence C5's count is exactly half C4's.}
\label{tab:refl}
\small
\setlength{\tabcolsep}{6pt}
\begin{tabular}{l r l r}
\toprule
Condition & Context instances & Unique strings & Success \\
\midrule
C3 memory$+$planning (no reflection) & -- & -- & 10\% \\
C4 deterministic (fair-share) & 47{,}040 & 3 (95.8\% top-2) & 70\% \\
\quad injection arm (string verbatim) & -- & 1 & 8/10 \\
\quad placebo arm (no content) & -- & 1 & 6/10 \\
\emph{C5 LLM reflection (exploratory)} & 23{,}520 & 1{,}190 (94.9\% norm) & 100\% \\
\bottomrule
\end{tabular}
\end{table}
\begin{figure}[t]
\centering
\IfFileExists{figures/fig_saturation.pdf}%
  {\includegraphics[width=0.68\textwidth]{figures/fig_saturation.pdf}}%
  {\figplaceholder{3.2cm}{Cumulative unique reflection strings vs.\ step, C4 and C5.}}
\caption{Cumulative unique reflection strings vs.\ step: C4 (deterministic, 3 unique strings, saturating by step 2) and C5 (LLM-generated, exploratory, 1{,}190 strings, still 94.9\% repeating the same norm).}
\label{fig:saturation}
\end{figure}






\section{Discussion and Threats to Validity}\label{sec:discussion}
\textbf{Implications for Web simulation practice.} (1)~Face validity alone is insufficient: believable agents reproduced a unanimous wrong consensus. (2)~Believability is a \emph{diagnostic signal}, not a gate on outcomes. (3)~Report the $(B,V)$ plane and discrepancy range alongside outcome metrics, since gap cases are what face validity most likely mis-certifies. (4)~Where a simulation exposes public information, test explicitly for herding. (5)~Architecture can help via unintended mechanisms (our reflection acted as a shared norm, not reasoning), so trace-level inspection~\cite{zhao2026mechanism,cemri2025fails} should accompany outcome metrics.




\subsection{Limitations}
\begin{itemize}[leftmargin=*,itemsep=1pt,topsep=2pt]
\item[] \textbf{Construct:} we measure task outcome reliability rather than simulation validity.
\item[] \textbf{Design:} believability covaries with architecture, so pooled associations mostly
reflect between-condition differences; the cumulative ladder does not identify individual module
effects. IC is a binary baseline-vs-enriched contrast (fine-grained statistics rest on CD)
confounded with option \emph{position}, read as cascade-like herding, not pure lock-in; a
no-flag-memory control shows memory still aggregates without the plurality hint.
\item[] \textbf{Statistics:} all primary trials share one 8-persona panel, so observations are not
independent and $p$-values are optimistic~\cite{wu2026boundary}; a four-panel mixed-effects
replication mitigates this (strong condition effect, negligible between-panel variance). We
foreground effect sizes given the multiple-testing surface; Type~E is a diagnostic label, not
independent evidence; and the presence-versus-quality distinction is observational, not causal.
\item[] \textbf{Instrument:} the composite uses a priori weights, supported only by a small
three-rater sanity check ($\rho=0.82$) that clears the reliability gate on no dimension. The
macro-overlapping cooperation term can be removed without changing the association, and
categorical results are instrument-independent. Planning and naturalness do share the generator's
judge model, though, a circularity a cross-provider judge would address.
\item[] \textbf{Generality:} one architecture~\cite{park2023generative} and one generator on two
tasks, so cross-architecture and cross-model replication is future work.
\end{itemize}


\section{Conclusion}\label{sec:conclusion}
We asked whether individual believability in LLM multi-agent simulations implies reliable collective outcomes; it does not. Across 120 primary and 20 exploratory trials, a positive but fragile association ($\rho=0.55$ in CD and $\rho=0.75$ in IC) is largely attributable to capability presence and disappears or reverses under not-applicable scoring. The clearest evidence is a public-voting task where believable agents converge unanimously on a wrong answer, confounded with option position, so we read it as cascade-like herding instead of pure lock-in; the coordination gain associated with deterministic reflection is substantially recovered by a content-free placebo, so it is not clearly attributable to the reflection content. Believability is thus a usable diagnostic signal, but reliable outcomes must be established separately. Future work should examine hidden-tally/vote-only IC controls, a direct aggregation-aid condition, full multi-rater validation, and cross-architecture/cross-model replication.

\subsubsection*{Acknowledgement.} We also acknowledge the use of AI-assisted language tools for
editorial refinement during the preparation of this manuscript.

\begin{thebibliography}{99}\scriptsize\linespread{0.92}\selectfont\setlength{\itemsep}{0pt}\setlength{\parsep}{0pt}\setlength{\parskip}{0pt}
\bibitem{park2023generative}
Park, J.S., O'Brien, J.C., Cai, C.J., et al.: Generative agents: interactive simulacra of human behavior. In: 36th Annual ACM Symp.\ User Interface Software and Technology (UIST '23), pp.\ 2:1--2:22. ACM, New York (2023). doi: 10.1145/3586183.3606763
\bibitem{gao2024survey}
Gao, C., et al.: Large language models empowered agent-based modeling and simulation: a survey and perspectives. Humanit. Soc. Sci. Commun. \textbf{11} (2024). doi: 10.1057/s41599-024-03611-3
\bibitem{ghaffarzadegan2024tutorial}
Ghaffarzadegan, N., Majumdar, A., Williams, R., Hosseinichimeh, N.: Generative agent-based modeling: an introduction and tutorial. System Dynamics Review \textbf{40}(4) (2024). doi: 10.1002/sdr.1761
\bibitem{yang2024oasis}
Yang, Z., et al.: OASIS: open agent social interaction simulations with one million agents. arXiv:2411.11581 (2024)
\bibitem{piao2025agentsociety}
Piao, J., et al.: AgentSociety: large-scale simulation of LLM-driven generative agents advances understanding of human behaviors and society. arXiv:2502.08691 (2025)
\bibitem{windrum2007empirical}
Windrum, P., Fagiolo, G., Moneta, A.: Empirical validation of agent-based models: alternatives and prospects. J. Artif. Soc. Social Simul. \textbf{10}(2) (2007)
\bibitem{moss2005towards}
Moss, S., Edmonds, B.: Towards good social science. J. Artif. Soc. Social Simul. \textbf{8}(4) (2005)
\bibitem{bikhchandani1992theory}
Bikhchandani, S., Hirshleifer, D., Welch, I.: A theory of fads, fashion, custom, and cultural change as informational cascades. J. Political Economy \textbf{100}(5), 992--1026 (1992). doi: 10.1086/261849
\bibitem{cho2025herd}
Cho, Y., Guntuku, S.C., Ungar, L.: Herd behavior: investigating peer influence in LLM-based multi-agent systems. arXiv:2505.21588 (2025)
\bibitem{jain2025social}
Jain, A., Krishnamurthy, V.: Interacting large language model agents: interpretable models and social learning. arXiv:2411.01271 (2025)
\bibitem{zhong2025conformity}
Zhong, H., et al.: Disentangling the drivers of LLM social conformity. arXiv:2508.14918 (2025)
\bibitem{zheng2023judging}
Zheng, L., et al.: Judging LLM-as-a-judge with MT-Bench and Chatbot Arena. In: Advances in Neural Information Processing Systems 36 (NeurIPS 2023), Datasets and Benchmarks Track. arXiv:2306.05685
\bibitem{aharon2026tacit}
Aharon, I., La Malfa, E., Wooldridge, M., et al.: Tacit coordination of large language models. arXiv:2601.22184 (2026)
\bibitem{epstein1996growing}
Epstein, J.M., Axtell, R.: Growing Artificial Societies: Social Science from the Bottom Up. Brookings Institution Press, Washington, DC (1996)
\bibitem{sargent2010verification}
Sargent, R.G.: Verification and validation of simulation models. In: 2010 Winter Simulation Conf.\ (WSC), pp.\ 166--183. IEEE, Piscataway (2010). doi: 10.1109/WSC.2010.5679166
\bibitem{rand2011rigor}
Rand, W., Rust, R.T.: Agent-based modeling in marketing: guidelines for rigor. Int. J. Research in Marketing \textbf{28}(3), 181--193 (2011). doi: 10.1016/j.ijresmar.2011.04.002
\bibitem{augusiak2014evaludation}
Augusiak, J., Van den Brink, P.J., Grimm, V.: Merging validation and evaluation of ecological models to `evaludation'. Ecol. Modelling \textbf{280}, 117--128 (2014). doi: 10.1016/j.ecolmodel.2013.11.009
\bibitem{zhao2026mechanism}
Zhao, P., Pham, D.H., Vincent, N.: Mechanism plausibility in generative agent-based modeling. In: 2026 ACM Conf.\ Fairness, Accountability, and Transparency (FAccT '26), pp.\ 5915--5940. ACM, New York (2026). doi: 10.1145/3805689.3812388
\bibitem{grimm2006odd}
Grimm, V., Berger, U., Bastiansen, F., et al.: A standard protocol for describing individual-based and agent-based models. Ecol. Modelling \textbf{198}(1--2), 115--126 (2006). doi: 10.1016/j.ecolmodel.2006.04.023
\bibitem{grimm2020odd}
Grimm, V., Railsback, S.F., Vincenot, C.E., et al.: The ODD protocol for describing agent-based and other simulation models: a second update. J. Artif. Soc. Social Simul. \textbf{23}(2), 7 (2020). doi: 10.18564/jasss.4259
\bibitem{grimm2005pom}
Grimm, V., Revilla, E., Berger, U., et al.: Pattern-oriented modeling of agent-based complex systems: lessons from ecology. Science \textbf{310}(5750), 987--991 (2005). doi: 10.1126/science.1116681
\bibitem{kluegl2008validation}
Kl\"ugl, F.: A validation methodology for agent-based simulations. In: 2008 ACM Symp.\ Applied Computing (SAC '08), pp.\ 39--43. ACM, New York (2008). doi: 10.1145/1363686.1363696
\bibitem{larooij2025critical}
Larooij, M., T{\"o}rnberg, P.: Validation is the central challenge for generative social simulation: a critical review of LLMs in agent-based modeling. Artif. Intell. Review \textbf{59}(1), article 15 (2026). doi: 10.1007/s10462-025-11412-6
\bibitem{zeng2026uncanny}
Zeng, Y., Brown, C., Rounsevell, M.: Too human to model: the uncanny valley of large language models in simulating human systems. npj Complexity \textbf{3}, article 13 (2026). doi: 10.1038/s44260-026-00075-1
\bibitem{wu2026boundary}
Wu, Z., et al.: LLM-based social simulations require a boundary. arXiv:2506.19806 (2025)
\bibitem{adornetto2025overview}
Adornetto, C., et al.: Generative agents in agent-based modeling: overview, validation, and emerging challenges. IEEE Trans. Artif. Intell. \textbf{6}(12), 3165--3183 (2025). doi: 10.1109/TAI.2025.3566362
\bibitem{cross2025validating}
Cross, L., Haber, N., Yamins, D.L.K.: Validating generative agent-based models of social norm enforcement: from replication to novel predictions. arXiv:2507.22049 (2025)
\bibitem{munker2026benchmark}
M{\"u}nker, S., Schwager, N., Rettinger, A.: Don't trust generative agents to mimic communication on social networks unless you benchmarked their empirical realism. In: 19th Conf.\ European Chapter Assoc.\ Computational Linguistics (EACL 2026), Vol.\ 1: Long Papers, pp.\ 1141--1151. ACL, Rabat (2026). doi: 10.18653/v1/2026.eacl-long.51
\bibitem{tomasevic2026operational}
Toma\v{s}evi\'c, A., et al.: Towards operational validation of LLM-agent social simulations: a replicated study of a Reddit-like technology forum. arXiv:2508.21740 (2026)
\bibitem{li2024agentbench}
Liu, X., et al.: AgentBench: evaluating LLMs as agents. In: International Conf.\ Learning Representations (ICLR 2024). arXiv:2308.03688
\bibitem{mohammadi2025evaluation}
Mohammadi, M., Li, Y., Lo, J., Yip, W.: Evaluation and benchmarking of LLM agents: a survey. In: 31st ACM SIGKDD Conf.\ Knowledge Discovery and Data Mining (KDD '25). ACM, New York (2025). doi: 10.1145/3711896.3736570
\bibitem{zhang2026silo}
Zhang, Y., et al.: Silo-Bench: a scalable environment for evaluating distributed coordination in multi-agent LLM systems. arXiv:2603.01045 (2026)
\bibitem{kulshreshtha2026defection}
Kulshreshtha, D., et al.: The subtle art of defection: understanding uncooperative behaviors in LLM-based multi-agent systems. arXiv:2511.15862 (2026)
\bibitem{hishiki2026memory}
Hishiki, T., Arita, T., Suzuki, R.: How memory can affect collective and cooperative behaviors in an LLM-based social particle swarm. arXiv:2604.12250 (2026)
\bibitem{galanti2020explainable}
Galanti, R., Coma-Puig, B., de Leoni, M., et al.: Explainable predictive process monitoring. In: 2nd Int.\ Conf.\ Process Mining (ICPM 2020), pp.\ 1--8. IEEE, Piscataway (2020). doi: 10.1109/ICPM49681.2020.00012
\bibitem{cemri2025fails}
Cemri, M., Pan, M.Z., Yang, S., et al.: Why do multi-agent LLM systems fail? arXiv:2503.13657 (2025)
\bibitem{yehudai2026clear}
Yehudai, A., et al.: Agentic CLEAR: automating multi-level evaluation of LLM agents. arXiv:2605.22608 (2026)
\bibitem{pan2025agentrewardbench}
L\`u, X.H., Kazemnejad, A., Meade, N., et al.: AgentRewardBench: evaluating automatic evaluations of web agent trajectories. arXiv:2504.08942 (2025)
\end{thebibliography}
\end{document}
